\documentclass[a4paper,12pt]{article}
\usepackage[utf8]{inputenc}
\usepackage[T1]{fontenc}
\usepackage[margin=2.5cm]{geometry}
\usepackage{ctex} % 中文支持
\usepackage{xcolor}
\usepackage{titlesec}
\usepackage{setspace}

% -----------------------------------------------------------
% 格式设置区：Level 3 纯文本墙 (Wall of Text)
% -----------------------------------------------------------

% 仅保留一级标题，格式简洁
\titleformat{\section}{\Large\bfseries\color{blue!40!black}}{\thesection}{1em}{}

% 设置行间距，在密集文本中保持一定的阅读舒适度
\onehalfspacing
\setlength{\parskip}{1em}
\setlength{\parindent}{2em}

% -----------------------------------------------------------
% 文档元数据
% -----------------------------------------------------------
\title{\textbf{2024年度环境、社会及管治（ESG）报告：可持续发展管理与绩效}}
\author{本公司董事会办公室}
\date{2025年3月}

\begin{document}

\maketitle

\section{前言与环境管理体系}

作为国内领先的化工行业上市公司，本公司始终将可持续发展作为企业运营的核心理念，严格遵守《环境保护法》及相关合规要求，致力于构建对环境友好的生产体系。2024年，我们建立并运行了完善的环境数据监测与管理机制，确保向利益相关方披露真实、准确的环境绩效数据。根据年度监测结果，本公司全年的温室气体排放总量控制在39,107.65吨CO2e，排放强度为4.58吨/百万元营收，这一数据体现了我们在低碳转型方面的努力。在废弃物管理合规方面，我们严格管控生产过程，全年产生一般工业固废170.48吨；对于环境风险较高的危险废物，我们执行了最严格的处置标准，处置率达到100\%，确保所有危废均得到合规的安全处理。在资源利用方面，我们坚持节约优先的原则，全年新鲜水总用量为149.55万吨，包装材料使用总量为7,677吨。为了支撑上述环境目标的实现，公司在2024年投入了专项环保资金457.93万元，并依法缴纳环保税9.6万元，切实履行了企业的环境责任。

为了进一步降低运营过程中的碳足迹，我们在生产和物流两个核心环节实施了具体的改进措施。在生产端，我们的核心生产基地响应国家清洁生产的号召，完成了清洁能源热能替代项目。这项工程彻底淘汰了老旧的燃煤锅炉，引入了清洁能源热源，就像把家里烧煤的炉子换成了清洁的天然气系统一样，该举措直接推动该基地实现了生产环节废气的“零排放”。同时，我们旗下的化工厂积极推行循环经济模式，实施了冷凝水回收利用项目，将生产蒸汽产生的热净水回收循环使用，既节约了水资源，又回收了热能。在物流端，针对西北某大型露天煤矿项目的运输排放问题，我们推进了车辆电气化战略。公司引入了21台纯电动矿卡，利用能量回收技术大幅降低能耗；补充了36台油电增程式矿卡作为过渡方案；并前瞻性地投入了12台无人驾驶电动矿卡，这些智能设备实现了运行过程的零排放和低噪音，最大限度地减少了对矿区生态环境的干扰。

\section{研发创新与技术赋能}

科技创新不仅是公司发展的动力，更是我们解决行业安全痛点、推动产业升级的责任所在。2024年，公司持续加大研发投入，金额达到4.19亿元，占营业收入的比例为4.91\%。我们组建了一支由1,064名专业人员构成的研发团队，并取得了丰硕的知识产权成果，全年新增授权专利126项，累计拥有专利596项。我们重点推广了三项具有行业示范意义的技术，旨在通过技术进步提升本质安全水平。首先是“现场混装水胶炸药工艺”，这项达到国际先进水平的技术改变了传统炸药的运输模式，实行原材料分开运输、现场即时混合。这不仅消除了运输途中的爆炸风险，还能根据岩石条件灵活调节威力，实现了安全与效益的双赢。其次是国内首创的“二氧化碳相变膨胀激发管”，该技术利用液态二氧化碳物理膨胀原理破岩，全过程无明火，如同给高瓦斯矿井穿上了“防爆铠甲”，极大提升了特殊环境下的作业安全性。最后是“智能矿山系统”，我们集成了5G、边缘计算与机器人技术，实现了5G钻机、自动验孔机器人与无人矿卡的协同作业，将工人从危险的一线环境中解放出来，真正体现了“科技以人为本”的发展理念。

\section{社会责任与员工关怀}

我们深知，员工的幸福和社区的繁荣是企业长青的基石。在员工权益保障方面，截至2024年底，公司员工总数为7,399人，劳动合同签订率保持在100\%。我们将安全生产视为对员工最庄严的承诺，建立了“1+14+N”的应急预案体系，并投入真金白银构筑安全防线，包括投入355.99万元购买安全生产责任险，以及投入892.5万元用于工伤保险。我们的安全管理体系在海外项目中也得到了有效验证。在纳米比亚矿服项目，面对跨文化管理的挑战，我们通过推行“晨会制度”和“手指口述”安全确认法等标准化管理手段，实现了连续130万小时无损工事件的优异成绩，树立了中国企业在海外合规运营的标杆。

作为央企控股企业，我们在履行社会责任方面始终冲锋在前。2024年，公司对外捐赠及公益总投入达到1,219万元，其中1,125万元专项用于支持乡村振兴，助力改善欠发达地区的基础设施和民生水平。在应对突发公共事件时，我们也展现了应有的担当。例如，在湖南某地发生堤坝决口险情时，我们迅速响应政府号召，调集专业矿山团队支援。依托精湛的爆破技术，工程师们采用特殊的“松动爆破”方案，在极短时间内开采出8.4万吨急需的石料，为抗洪抢险争取了宝贵时间。这一行动充分证明了我们将专业能力转化为社会价值的决心。

\section{公司治理与商业道德}

完善的公司治理结构和高标准的商业道德是保护投资者利益的前提。作为上级集团的重要成员，我们致力于构建透明、高效的治理机制。我们的董事会采用了科学的“3+3+3”结构，即3名股东董事、3名独立董事与3名执行董事，这种平衡的架构有效保障了决策的独立性与制衡性，避免了决策风险。凭借规范的治理实践，公司在证券交易所的信息披露考评中获得了“A级”评价。在股东回报方面，我们坚持与投资者共享发展成果，2024年共派发现金红利2.54亿元，每10股派发2.05元，现金分红比例达到可分配利润的40.12\%。在供应链合规管理方面，我们高度重视廉洁风险防控。针对合作的1,548家供应商，我们实施了严格的准入机制，要求所有供应商必须签署《廉洁协议》。截至目前，该协议的签署率达到100\%，这一举措有效防范了商业贿赂风险，确保了采购过程的公开、透明，为构建健康的行业生态做出了积极贡献。

\end{document}